\section{Conclusion}\label{sec:conclusion}

The finite-window fiber spectrum admits two exact descriptions.  The affine
permutation \(\pi_m\) identifies each multiplicity with a coefficient of a
finite Fibonacci subset-sum polynomial and with a Fibonacci-lag difference of
the indexed partition function.  The interval identity identifies the same
fiber list with ordinary partition values on two consecutive Fibonacci
layers.  The latter correspondence transfers the published extremal
classification, after an explicit convention change, and yields
\[
 D_{2k}=F_{k+2},\qquad D_{2k+1}=2F_{k+1},\qquad
 D_m=D_{m-2}+D_{m-4}\quad(m\ge6),
\]
together with the complete maximizing-residue list.

For positive integer orders, the adjacent-window sandwich transfers Sanna's
partition power-sum constants to the finite-window moments.  Sanna's endpoint
law \(\lambda_q^{1/q}\to\varphi^{1/2}\) also follows here from the maximum-norm
squeeze; the separate consequence of that squeeze is the diagonal limit for
\(q_m\to\infty\), and indeed for arbitrary real \(t_m\to\infty\).  The
quadratic moment has a direct cubic recurrence that does not use finite-state
normalization.

For negative real tilts, Weinstein's orbit counts give the frozen half-line
\[
 P(t)=0\qquad(t\le-\sigma_0),
 \qquad
 \frac{\zeta(\sigma_0-1)}{\zeta(\sigma_0)}=2.
\]
The uniform layer-count lemma treats the two unit conventions, the levels
\(1\) and \(2\), shared endpoints, and early layers explicitly.  Analytic
perturbation theory then rules out a representation of all real pressure
limits by the positive simple isolated eigenvalue of an analytic compact
operator family satisfying the stated local gap condition.

Existence and regularity of \(P(t)\) on
\(( -\sigma_0,\infty)\) remain open away from the positive integers.  The full
real-tilt large deviation principle is therefore only the conditional
corollary in \Cref{thm:conditional-full-ldp}.  No transducer-based collision
kernel, algebraicity statement for \(\lambda_q\), arithmetic-window recurrence
polynomial, exact-kernel certificate, Galois group, or Chebotarev density is
claimed in this paper.
